\section{问题二的建模与求解}

\subsection{模型基本设定}

问题二沿用问题一的时间划分、储能状态转移及充放电约束，并进一步引入日期下标 $d$ 和紧急购电量 $e_{d,t}$。由于问题二未要求储能设备每日0:00与24:00的储电量相同，储电量在相邻日期间连续传递：
\begin{equation}
E_{d+1,0}=E_{d,144},
\qquad
E_{2025\text{-}01\text{-}01,0}=6000\ \mathrm{kWh}.
\label{eq:p2_crossday}
\end{equation}

\subsection{完全信息基准模型}

为衡量日前信息不完全带来的额外购电成本，我们来构造完全信息基准。将附件2中的实际负载与实际光伏发电量视为优化前已知信息，在微网购电与储能基础约束下求解完全信息条件下的最优调度结果。设$x_{d,t}$和$e_{d,t}$分别表示第$d$天第$t$时段的正常购电量和紧急购电量，$c_t$为对应时段的外部电网电价，则目标函数为
\begin{equation}
\min C_A=
\sum_{d}\sum_{t}
\left(c_t x_{d,t}+5c_t e_{d,t}\right).
\label{eq:p2_oracle}
\end{equation}

完全信息条件下，各时段实际负载和光伏发电量均已知，且正常购电价格低于紧急购电价格，因此模型能够提前通过正常购电与储能调度满足供电需求。此时正常购电量为$20\,182\,315.798\ \mathrm{kWh}$，紧急购电量为0。将该结果作为完全信息条件下的理论基准，用于衡量日前信息有限产生的额外购电成本。


\subsection{基于历史信息的日前购电策略}

\subsubsection{基于前一日数据的供需估计}

完全信息基准依赖目标日实际负载和光伏发电量，无法直接用于日前决策。考虑相邻日期供需曲线具有一定连续性，在目标日$d$的0:00，首先采用前一日实际数据估计当天各时段的负载和光伏发电量：
\begin{equation}
\hat L_{d,t}=L_{d-1,t},\qquad
\hat R_{d,t}=R_{d-1,t}.
\label{eq:p2_yesterday}
\end{equation}

将$\hat L_{d,t}$和$\hat R_{d,t}$代入问题一的储能调度模型，确定目标日各时段的计划购电量及储能充放电方案。日前计划确定后，根据当天实际负载和实际光伏发电量进行结算；当计划供电无法满足实际需求时，缺口部分通过紧急购电补足：
\begin{equation}
e_{d,t}
=
\max\left\{
0,\,
L_{d,t}^{\mathrm{act}}+u_{d,t}
-x_{d,t}-R_{d,t}^{\mathrm{act}}-v_{d,t}
\right\}.
\label{eq:p2_emergency}
\end{equation}


\subsubsection{固定安全冗余策略}

仅采用前一日数据进行预测时，日间供需变化可能造成计划购电不足。
考虑紧急购电价格为交易时刻电价的5倍，在基础计划购电量上，
对每个10 min时段额外增加固定安全冗余量$\delta$：
\begin{equation}
x_{d,t}^{\mathrm{r}}
=
x_{d,t}^{\mathrm{base}}+\delta,
\qquad \delta\geq0.
\label{eq:p2_fixed_reserve}
\end{equation}

比较不同冗余水平下的正常购电费用、紧急购电费用及总购电费用后，
选取$\delta=70\ \mathrm{kWh}$作为固定安全冗余策略，
并与后续动态修正策略进行比较。


\subsection{动态修正购电策略}

负载与光伏发电量具有一定的星期规律和季节差异，统一采用前一日数据及固定安全
冗余难以充分反映不同日期的供需变化。因此，
进一步利用星期规律、季节特征和近期净负荷偏差修正预测结果，
并结合日初储电量动态确定安全冗余。


\subsubsection{历史规律预测}

小区负载可能受到不同星期生活节奏的影响。本文将星期一至星期日分别作为
七种星期类型，并将与目标日星期类型相同的历史日期定义为同星期历史日。

为选择历史预测误差较小的预测方式，以净负荷
\[
N_{d,t}=L_{d,t}-R_{d,t}
\]
作为评价对象，在目标日前最近7个已发生日期上进行滚动检验。
设$V_d$为目标日前最近7个已发生日期组成的验证集，
则第$k$种候选预测的平均绝对误差为
\begin{equation}
\mathrm{MAE}_k(d)
=
\frac{1}{|V_d|T}
\sum_{j\in V_d}\sum_{t=1}^{T}
\left|
N_{j,t}-\hat N^{(k)}_{j,t}
\right|.
\label{eq:p2_mae}
\end{equation}

比较不同候选方案的$\mathrm{MAE}$，选择误差较小者作为目标日基础预测；
误差相同时优先采用同星期历史数据。对统计期334个评价日逐日检验后，
均选择最近两个同星期历史日的平均曲线。

因此，选取目标日前最近两个相同星期类型的历史日，
取其平均负载和光伏发电量作为候选预测：
\begin{equation}[H]
\left(\hat L^{W}_{d,t},\hat R^{W}_{d,t}\right)
=
\frac{1}{|H_d^W|}
\sum_{h\in H_d^W}
\left(L_{h,t},R_{h,t}\right),
\label{eq:p2_week_candidate}
\end{equation}
其中，$H_d^W$表示目标日前最近两个与目标日星期类型相同的历史日期集合。


\subsubsection{日期与季节性修正}

题目表3给出的四个指定日期分别对应春分、夏至、秋分和冬至，具有明显的季节代表性。

\begin{figure}[H]
    \centering
    \includegraphics[width=0.8\linewidth]{图04_典型日期小区负载与光伏实际功率.png}
    \caption{典型日期小区负载与光伏实际功率}
    \label{fig:04}
\end{figure}

根据附件2中对应日期的实际数据，绘制典型日期的小区负载与光伏发电功率，如图~\ref{fig:04}所示。

不同季节的光伏发电时段和小区整体用电水平存在差异。一方面，太阳高度角随季节变化，不同季节发电峰值出现时刻和日内曲线形状不尽相同;另一方面，季节变化可能影响整体用电水平。例如夏季制冷、冬季取暖及不同季节的生活作息，都可能使全天负荷总量和峰谷特征发生变化。

为保留历史曲线基本特征并反映日期变化，根据日期对候选光伏发电曲线的
时间分布进行修正，并以25\%的权重与原曲线融合，同时保持候选日光伏发电
总量不变。对于小区负载，则根据历史负载总量随日期的变化规律修正候选曲线
整体水平，并以50\%的权重与原预测结果融合。




\subsubsection{近期偏差修正}

为进一步反映近期供需变化，根据目标日前最近7天的净负荷预测残差构造
逐时段偏差修正量$b_{d,t}$。对近期误差赋予较高权重，并对异常偏差进行
截断处理，修正后的日前预测为
\begin{equation}
\begin{aligned}
\hat L_{d,t}^{\,\mathrm{base}}
&=
\hat L_{d,t}^{\,\mathrm{his}}
+\lambda b_{d,t},\\
\hat R_{d,t}^{\,\mathrm{base}}
&=
\hat R_{d,t}^{\,\mathrm{his}},
\end{aligned}
\qquad \lambda=0.75.
\label{eq:p2_bias_correction}
\end{equation}

当近期实际净负荷持续高于基准预测值时，通过提高日前负荷需求估计，
降低计划购电不足的风险。


\subsubsection{动态安全冗余}

考虑不同日期的预测误差及日初储电量存在差异，将固定安全冗余进一步扩展为
随日期变化的动态安全冗余量。每日从候选集合
\begin{equation}
\delta_d
\in
\{0,10,20,\ldots,140\}
\quad \mathrm{kWh}
\label{eq:p2_delta_set}
\end{equation}
中确定$\delta_d$。

其中，$\delta_d$表示目标日每个10 min时段增加的安全冗余电量，
同一日期各时段采用相同数值。综合历史净负荷偏差、正常购电成本及日初储电量，
在预测误差较大或储能余量较低时提高安全冗余。

将动态安全冗余直接加入日前负荷需求并参与储能联合优化：
\begin{equation}
\hat L_{d,t}^{\,\mathrm{MILP}}
=
\hat L_{d,t}^{\,\mathrm{base}}
+\delta_d.
\label{eq:p2_dynamic_reserve}
\end{equation}

由此形成Dynamic日前购电策略，使安全冗余与计划购电、储能充放电共同参与优化。


\subsubsection{实际结算与盈余划分}

日前计划确定后，采用目标日实际负载和实际光伏发电量进行结算。
计划供给不足时按照前述紧急购电规则补足；实际供给超过负载及储能需求时，
定义剩余供给量为
\begin{equation}
s_{d,t}
=
\max\left\{
0,\,
x_{d,t}+R_{d,t}^{\mathrm{actual}}+v_{d,t}
-L_{d,t}^{\mathrm{actual}}-u_{d,t}
\right\}.
\label{eq:p2_surplus}
\end{equation}

进一步将剩余供给划分为光伏弃光量和外部电网购电溢余量：
\begin{equation}
w_{d,t}^{PV}
=
\min\left\{
s_{d,t},R_{d,t}^{\mathrm{actual}}
\right\},
\label{eq:p2_pv_curtailment}
\end{equation}

\begin{equation}
g_{d,t}^{\mathrm{spill}}
=
\max\left\{
s_{d,t}-R_{d,t}^{\mathrm{actual}},0
\right\}.
\label{eq:p2_grid_spill}
\end{equation}

其中，$w_{d,t}^{PV}$表示实际光伏发电量中未被负载及储能利用的部分，
$g_{d,t}^{\mathrm{spill}}$表示外部电网购电超过实际负载及储能需求后形成的
剩余电量。上述划分用于统计不同购电策略下的光伏消纳与计划购电过量情况。


\subsection{不同购电策略的结果比较}

为比较不同信息条件与安全冗余策略的购电效果，分别计算完全信息基准、
Yesterday-MILP、Yesterday-MILP+70和Dynamic四种策略在统计期内的运行结果。
除完全信息基准外，其余策略均仅使用决策时点已经获得的信息制定日前计划，
并根据目标日实际负载和实际光伏发电量进行结算，结果见表~\ref{tab:p2_compare}。

\begin{table}[H]
\centering
\caption{不同购电策略的统计期运行结果}
\label{tab:p2_compare}
\small
\resizebox{\textwidth}{!}{
\begin{tabular}{llrrrrrrr}
\toprule
模型 & 利用信息 & 计划购电量 & 计划购电费用 & 实际紧急购电量 &
紧急购电费用 & 实际结算总购电费用 & 弃光量 & 外部电网购电溢余量 \\
\midrule
完全信息基准（A）
& 统计期实际信息
& 20,182,315.798
& 12,227,243.643
& 0.000
& 0.000
& 12,227,243.643
& 990,168.090
& 0.000 \\

Yesterday-MILP
& 前一日实际信息
& 20,139,144.857
& 12,217,149.502
& 2,588,654.457
& 10,606,597.790
& 22,823,747.293
& 2,457,440.902
& 1,071,499.546 \\

Yesterday-MILP+70
& 前一日实际信息
& 23,505,864.857
& 14,796,721.014
& 1,572,201.133
& 6,664,514.990
& 21,461,236.004
& 3,487,135.843
& 2,392,071.280 \\

Dynamic
& 历史信息及动态修正
& 22,556,016.294
& 13,861,172.313
& 318,742.269
& 1,208,683.712
& 15,069,856.025
& 2,325,197.192
& 1,301,651.115 \\
\bottomrule
\end{tabular}}
\begin{flushleft}
\footnotesize 注：金额单位为元，电量单位为kWh。
\end{flushleft}
\end{table}
Yesterday-MILP仅采用前一日实际曲线制定计划，统计期紧急购电量达到
$2,588,654.457$ kWh，实际结算总购电费用为$22,823,747.293$元。
增加70 kWh固定安全冗余后，紧急购电量下降至$1,572,201.133$ kWh，
总购电费用降至$21,461,236.004$元，但外部电网购电溢余量增加至
$2,392,071.280$ kWh。

Dynamic进一步修正日前供需预测，并将动态安全冗余纳入储能联合优化。
统计期紧急购电量降至$318,742.269$ kWh，实际结算总购电费用为
$15,069,856.025$元。相较Yesterday-MILP+70，紧急购电量减少约$79.7\%$，
总购电费用减少约$639.1$万元，因此选取Dynamic作为问题二的最终购电策略。

与完全信息基准相比，Dynamic实际结算总购电费用增加
$2,842,612.382$元，约高$23.25\%$，该差额反映有限信息条件下预测偏差等因素
产生的额外购电成本。

\begin{figure}[H]
    \centering
    \includegraphics[width=0.75\linewidth]{图05_不同购电策略间的总费用与紧急购电量对比.png}
    \caption{不同购电策略间的总费用与紧急购电量对比}
    \label{fig:05}
\end{figure}


\subsection{指定日期的策略比较}

进一步考察题目指定的2025年3月20日、6月21日、9月23日和12月21日，
比较完全信息基准与Dynamic的日前计划及实际结算结果，
如表~\ref{tab:p2_selected_days}所示。其中，$\delta_d$表示Dynamic在目标日
采用的安全冗余量，日初和日末储电量由跨日连续运行结果确定。

\begin{table}[H]
\centering
\caption{指定日期不同购电策略的运行结果}
\label{tab:p2_selected_days}
\small
\resizebox{\textwidth}{!}{
\begin{tabular}{llrrrrrrr}
\toprule
日期 & 模型 &
计划购电费用/元 &
紧急购电量/kWh &
紧急购电费用/元 &
总购电费用/元 &
$\delta_d$/每10 min/kWh &
日初储电量/kWh &
日末储电量/kWh \\
\midrule

2025-03-20 & 完全信息基准（A）
& 37,960.409 & 0.000 & 0.000 & 37,960.409
& -- & 8,550.000 & 4,352.866 \\

2025-03-20 & Dynamic
& 43,421.767 & 302.639 & 1,258.158 & 44,679.925
& 30 & 1,200.000 & 1,200.000 \\

2025-06-21 & 完全信息基准（A）
& 20,470.623 & 0.000 & 0.000 & 20,470.623
& -- & 2,361.060 & 8,550.000 \\

2025-06-21 & Dynamic
& 20,611.914 & 362.215 & 1,308.939 & 21,920.853
& 20 & 1,200.000 & 1,200.000 \\

2025-09-23 & 完全信息基准（A）
& 41,311.883 & 0.000 & 0.000 & 41,311.883
& -- & 8,550.000 & 8,550.000 \\

2025-09-23 & Dynamic
& 45,925.485 & 1,079.032 & 4,938.765 & 50,864.250
& 70 & 1,200.000 & 1,200.000 \\

2025-12-21 & 完全信息基准（A）
& 59,664.197 & 0.000 & 0.000 & 59,664.197
& -- & 8,550.000 & 8,550.000 \\

2025-12-21 & Dynamic
& 62,994.630 & 881.389 & 2,717.806 & 65,712.436
& 30 & 1,200.000 & 1,200.000 \\
\bottomrule
\end{tabular}}
\end{table}

Dynamic仅利用决策时点前已经获得的信息制定日前计划，因此实际负载和
光伏发电量与预测值之间的偏差仍可能引起紧急购电。
四个指定日期中，9月23日紧急购电量最高，为$1,079.032$ kWh，
对应每个10 min时段70 kWh的安全冗余；6月21日紧急购电量为
$362.215$ kWh，对应安全冗余为20 kWh。
不同日期采用的安全冗余量随历史预测偏差及储能状态进行调整。


\subsection{问题二购电策略与结果}

根据上述比较，采用\textbf{Dynamic}作为问题二的最终日前购电策略。\textbf{该策略基于目标日前历史数据预测当日负载与光伏发电量，并结合近期净负荷预测偏差和日初储电量动态设置安全冗余；在此基础上求解混合整数线性规划模型，确定各时段计划购电量及储能充放电方案，以降低预测偏差导致的紧急购电风险。}

表~\ref{tab:p2_purchase}展示题目指定时段的计划购电量以及
全天计划购电量和购电费用；表~\ref{tab:p2_storage}展示各4 h区间的储能
充放电量以及0:00、24:00储电量；表~\ref{tab:p2_emergency}展示紧急购电结果。


\begin{table}[H]
\centering
\caption{指定日期在指定时段的计划购电量、全天购电量与购电费用（Dynamic）}
\label{tab:p2_purchase}
\resizebox{\textwidth}{!}{
\begin{tabular}{cc cc cc cc}
\toprule
\multicolumn{2}{c}{2025-03-20} &
\multicolumn{2}{c}{2025-06-21} &
\multicolumn{2}{c}{2025-09-23} &
\multicolumn{2}{c}{2025-12-21} \\
\cmidrule(lr){1-2}\cmidrule(lr){3-4}
\cmidrule(lr){5-6}\cmidrule(lr){7-8}
时间段 & 购电量（kWh） &
时间段 & 购电量（kWh） &
时间段 & 购电量（kWh） &
时间段 & 购电量（kWh） \\
\midrule

10:00--10:10 & 0.0000 &
10:00--10:10 & 0.0000 &
10:00--10:10 & 0.0000 &
10:00--10:10 & 0.0000 \\

12:00--12:10 & 627.3723 &
12:00--12:10 & 0.0000 &
12:00--12:10 & 437.4294 &
12:00--12:10 & 907.5325 \\

14:00--14:10 & 0.0000 &
14:00--14:10 & 0.0000 &
14:00--14:10 & 0.0000 &
14:00--14:10 & 0.0000 \\

16:00--16:10 & 534.4301 &
16:00--16:10 & 215.4195 &
16:00--16:10 & 575.5673 &
16:00--16:10 & 856.2892 \\

18:00--18:10 & 729.2638 &
18:00--18:10 & 0.0000 &
18:00--18:10 & 847.3361 &
18:00--18:10 & 736.7735 \\

20:00--20:10 & 0.0000 &
20:00--20:10 & 0.0000 &
20:00--20:10 & 30.6555 &
20:00--20:10 & 0.0000 \\
\midrule

全天计划购电量（kWh） & 71402.7837 &
全天计划购电量（kWh） & 35799.6645 &
全天计划购电量（kWh） & 73017.3009 &
全天计划购电量（kWh） & 97961.2955 \\

全天计划购电费用（元） & 43,421.7669 &
全天计划购电费用（元） & 20,611.9143 &
全天计划购电费用（元） & 45,925.4846 &
全天计划购电费用（元） & 62,994.6298 \\
\bottomrule
\end{tabular}}
\end{table}

由表~\ref{tab:p2_purchase}可得，3月20日、6月21日、9月23日和12月21日的
全天计划购电量分别为$71,402.784$、$35,799.665$、$73,017.301$和
$97,961.296$ kWh，对应计划购电费用分别为$43,421.767$、$20,611.914$、
$45,925.485$和$62,994.630$元。题目指定时段的计划购电量均列于表中，
其余各10 min时段的完整计划购电结果见附件。


\begin{table}[H]
\centering
\caption{指定日期储能充放电量及日初、日末储电量（Dynamic）}
\label{tab:p2_storage}
\resizebox{\textwidth}{!}{
\begin{tabular}{ccc ccc ccc ccc}
\toprule
\multicolumn{3}{c}{2025-03-20} &
\multicolumn{3}{c}{2025-06-21} &
\multicolumn{3}{c}{2025-09-23} &
\multicolumn{3}{c}{2025-12-21} \\
\cmidrule(lr){1-3}\cmidrule(lr){4-6}
\cmidrule(lr){7-9}\cmidrule(lr){10-12}
时间段 & 充电量 & 放电量 &
时间段 & 充电量 & 放电量 &
时间段 & 充电量 & 放电量 &
时间段 & 充电量 & 放电量 \\
 & （kWh） & （kWh） &
 & （kWh） & （kWh） &
 & （kWh） & （kWh） &
 & （kWh） & （kWh） \\
\midrule

0:00--4:00 & 9000.0000 & 0.0000 &
0:00--4:00 & 1776.0011 & 0.0000 &
0:00--4:00 & 9000.0000 & 0.0000 &
0:00--4:00 & 9000.0000 & 0.0000 \\

4:00--8:00 & 1666.6667 & 5585.8133 &
4:00--8:00 & 833.3333 & 2113.5609 &
4:00--8:00 & 1666.6667 & 6630.5465 &
4:00--8:00 & 2500.0000 & 1508.3333 \\

8:00--12:00 & 5876.8755 & 3054.1867 &
8:00--12:00 & 4938.1956 & 0.0000 &
8:00--12:00 & 5495.9677 & 2009.4535 &
8:00--12:00 & 5833.3333 & 7806.6667 \\

12:00--16:00 & 5350.5516 & 454.2160 &
12:00--16:00 & 5728.4710 & 0.0000 &
12:00--16:00 & 5587.1587 & 337.3324 &
12:00--16:00 & 8051.6126 & 2606.8062 \\

16:00--20:00 & 0.0000 & 5627.6462 &
16:00--20:00 & 0.0000 & 6008.6600 &
16:00--20:00 & 0.0000 & 5833.3333 &
16:00--20:00 & 0.0000 & 5559.6742 \\

20:00--24:00 & 0.0000 & 3012.3538 &
20:00--24:00 & 0.0000 & 2631.3400 &
20:00--24:00 & 0.0000 & 2806.6667 &
20:00--24:00 & 0.0000 & 3080.3258 \\
\midrule

0:00储电量 & 1200.0000 & -- &
0:00储电量 & 1200.0000 & -- &
0:00储电量 & 1200.0000 & -- &
0:00储电量 & 1200.0000 & -- \\

24:00储电量 & 1200.0000 & -- &
24:00储电量 & 1200.0000 & -- &
24:00储电量 & 1200.0000 & -- &
24:00储电量 & 1200.0000 & -- \\
\bottomrule
\end{tabular}}
\end{table}

表~\ref{tab:p2_storage}给出了四个指定日期各4 h区间的储能充、放电量。
四个日期0:00与24:00储电量均为$1200.000$ kWh。
问题二不要求每日结束时储电量恢复至日初水平，储能状态按照跨日连续运行方式
传递；各10 min时段的完整储能调度结果见附件。


由于紧急购电涉及全天144个10 min时段，正文仅展示四个指定日期的前若干时段，
并汇总各日紧急购电总量；完整非零紧急购电区间及对应购电量见附录（待补充）。

\begin{table}[H]
\centering
\caption{微电网在指定日期的紧急购电量（Dynamic）}
\label{tab:p2_emergency}
\resizebox{\textwidth}{!}{%
\begin{tabular}{cccccccc}
\toprule
\multicolumn{2}{c}{2025-03-20} &
\multicolumn{2}{c}{2025-06-21} &
\multicolumn{2}{c}{2025-09-23} &
\multicolumn{2}{c}{2025-12-21} \\
\cmidrule(lr){1-2}
\cmidrule(lr){3-4}
\cmidrule(lr){5-6}
\cmidrule(lr){7-8}
时间段 & 购电量（kWh） &
时间段 & 购电量（kWh） &
时间段 & 购电量（kWh） &
时间段 & 购电量（kWh） \\
\midrule

0:00--0:10 & 0.0000 &
0:00--0:10 & 0.0000 &
0:00--0:10 & 0.0000 &
0:00--0:10 & 0.0000 \\

0:10--0:20 & 0.0000 &
0:10--0:20 & 0.0000 &
0:10--0:20 & 0.0000 &
0:10--0:20 & 0.0000 \\

0:20--0:30 & 0.0000 &
0:20--0:30 & 34.9023 &
0:20--0:30 & 0.0000 &
0:20--0:30 & 0.0000 \\

0:30--0:40 & 0.0000 &
0:30--0:40 & 0.0000 &
0:30--0:40 & 0.0000 &
0:30--0:40 & 0.0000 \\

0:40--0:50 & 0.0000 &
0:40--0:50 & 0.0000 &
0:40--0:50 & 0.0000 &
0:40--0:50 & 0.0000 \\

0:50--1:00 & 0.0000 &
0:50--1:00 & 0.0000 &
0:50--1:00 & 0.0000 &
0:50--1:00 & 0.0000 \\

1:00--1:10 & 3.2300 &
1:00--1:10 & 0.0000 &
1:00--1:10 & 0.0000 &
1:00--1:10 & 0.0000 \\

1:10--1:20 & 21.2577 &
1:10--1:20 & 0.0000 &
1:10--1:20 & 0.0000 &
1:10--1:20 & 3.7376 \\

1:20--1:30 & 17.6671 &
1:20--1:30 & 38.2221 &
1:20--1:30 & 0.0000 &
1:20--1:30 & 0.0000 \\

$\cdots$ & $\cdots$ &
$\cdots$ & $\cdots$ &
$\cdots$ & $\cdots$ &
$\cdots$ & $\cdots$ \\
\midrule

\textbf{当日合计} & \textbf{302.6394} &
\textbf{当日合计} & \textbf{362.2149} &
\textbf{当日合计} & \textbf{1079.0316} &
\textbf{当日合计} & \textbf{881.3887} \\
\bottomrule
\end{tabular}%
}
\end{table}

四个指定日期的紧急购电总量分别为$302.639$、$362.215$、
$1,079.032$和$881.389$ kWh，其中9月23日最高，3月20日最低。
结合表~\ref{tab:p2_purchase}和表~\ref{tab:p2_storage}，
指定时段计划购电量、全天计划购电量与购电费用、
各4 h区间储能充放电量、日初与日末储电量以及紧急购电量均已给出；未在正文展开的逐时段计划购电、储能调度及紧急购电结果详见附件。